\begin{flushleft}
    Универзитет у Нишу \\

    Електронски факултет \\
    
    \vspace{5mm}
    
\end{flushleft}

    \noindent \textsc{\textbf{Класификација биомедицинских сигнала применом временско-фреквенцијске анализе на уређајима са орграниченим ресурсима}}
    
    \noindent \textsc{\textbf{Biomedical Signal Classification using Time-Frequency Analysis on Edge Devices}}\\
    
    
\begin{flushleft}

    Мастер рад\\
    
    Студијски програм: Електроника и микросистеми\\
    
    \vspace{5mm}
    
    Студент: Сандра Станковић, бр. инд. 1795\\
    
    Ментор: доц. др Стевица Цветковић
    \vspace{5mm}\\
    \end{flushleft}
    
    \noindent Задатак: \textit{Представити основне карактеристике фотоплетизмографског и акцелерометарског сигнала у контексту детекциjе стресних интервала. Развити методу за трансформациjу улазног сигнала у временско-фреквенциjски домен применом краткотрајне Фуриjеове трансформациjе (STFT). Пројектовати и имплементирати дубоку конволуциону неуронску мрежу коjа прецизно класификуjе стања стреса. Испитати утицаj различитих облика улазне STFT матрице на перформансе неуронске мреже. Оптимизовати модел за примену на уређаjима са ограниченим ресурсима и извршити поређење добиjених резултата са референтним резултатима из литературе.}

    
    \begin{flushleft}
    
    Датум пријаве рада: 11.09.2025.
    
    Датум предаје рада: 17.09.2025.
    
    Датум одбране рада: 19.09.2025. \\
    
    \vspace{5mm}
\end{flushleft}


\begin{flushright}
    Комисија за одбрану и оцену:
    
    \vspace{10mm}
    \begin{tabular}{@{}>{\raggedleft\arraybackslash}p{4in}@{}}
    \hrulefill \\
     1. Доц. др Стевица Цветковић, председник Комисије
    \end{tabular}

    \vspace{10mm}
    \begin{tabular}{@{}>{\raggedleft\arraybackslash}p{4in}@{}}
    \hrulefill \\
     2. Проф. др Горан Станчић, члан
    \end{tabular}

    \vspace{10mm}
    \begin{tabular}{@{}>{\raggedleft\arraybackslash}p{4in}@{}}
    \hrulefill \\
     3. Проф. др Саша В. Николић, члан
    \end{tabular}
\end{flushright}
